% Compact illustrations, using the supplied review packet's ordinary TikZ style.
\newcommand{\FunctionIllustration}{%
\begin{tikzpicture}[x=1cm,y=1cm,>=Stealth,font=\small]
\foreach \dx/\which in {0/1,6.6/2}{
 \begin{scope}[xshift=\dx cm]
  \node at (0,3.25) {$X$}; \node at (2.7,3.25) {$Y$};
 \end{scope}}
% Fixed node positions; only the subset appropriate to each operation is shaded.
\foreach \dx in {0,6.6}{
 \begin{scope}[xshift=\dx cm]
  \coordinate (a) at (0,2.35); \coordinate (b) at (0,1.65);
  \coordinate (c) at (0,.95); \coordinate (d) at (0,.25);
  \coordinate (p) at (2.7,2.35); \coordinate (q) at (2.7,1.3); \coordinate (r) at (2.7,.25);
  \draw[gray,rounded corners] (-.45,-.05) rectangle (.45,2.75);
  \draw[gray,rounded corners] (2.25,-.05) rectangle (3.15,2.75);
 \end{scope}}
% Image of U={a,c}.
\fill[blue!8,rounded corners] (-.3,2.12) rectangle (.3,2.58);
\fill[blue!8,rounded corners] (-.3,.72) rectangle (.3,1.18);
\fill[blue!8,rounded corners] (2.4,1.07) rectangle (3,2.58);
\foreach \name/\xx/\yy in {a/0/2.35,b/0/1.65,c/0/.95,d/0/.25,p/2.7/2.35,q/2.7/1.3,r/2.7/.25}
 \node (im\name) at (\xx,\yy) {$\name$};
\draw[->,blue] (ima)--(imp); \draw[->,gray] (imb)--(imp);
\draw[->,blue] (imc)--(imq); \draw[->,gray] (imd)--(imq);
\node[below] at (0,-.25) {$U=\{a,c\}$};
\node[below] at (2.7,-.25) {$f(U)=\{p,q\}$};
% Inverse image of V={q,r}.
\begin{scope}[xshift=6.6cm]
\fill[blue!8,rounded corners] (-.3,.02) rectangle (.3,1.18);
\fill[blue!8,rounded corners] (2.4,.02) rectangle (3,1.53);
\foreach \name/\xx/\yy in {a/0/2.35,b/0/1.65,c/0/.95,d/0/.25,p/2.7/2.35,q/2.7/1.3,r/2.7/.25}
 \node (pre\name) at (\xx,\yy) {$\name$};
\draw[->,gray] (prea)--(prep); \draw[->,gray] (preb)--(prep);
\draw[->,blue] (prec)--(preq); \draw[->,blue] (pred)--(preq);
\node[below] at (0,-.25) {$f^{-1}(V)=\{c,d\}$};
\node[below] at (2.7,-.25) {$V=\{q,r\}$};
\end{scope}
\end{tikzpicture}
\figcap{An image collects the outputs of the chosen inputs. An inverse image collects every input whose output is in the chosen set. No arrow ends at $r$, so $r$ contributes no input to $f^{-1}(V)$.}}

\newcommand{\ConjugateIllustration}{%
\begin{tikzpicture}[x=.6cm,y=.48cm,>=Stealth,font=\small]
\draw[->,gray] (-1,0)--(6.8,0) node[right,text=black] {real};
\draw[->,gray] (0,-5)--(0,5) node[above,text=black] {imaginary};
\draw[->,blue] (0,0)--(3,4) node[above right] {$z=3+4i$};
\draw[->,red] (0,0)--(3,-4) node[below right] {$\overline z=3-4i$};
\draw[dashed,gray] (3,4)--(3,-4);
\draw[dashed,gray] (0,4)--(3,4); \draw[dashed,gray] (0,-4)--(3,-4);
\node[below left] at (0,0) {$0$}; \node[below right] at (3,0) {$3$};
\node[left] at (0,4) {$4$}; \node[left] at (0,-4) {$-4$};
\end{tikzpicture}
\figcap{Conjugation changes the sign of the imaginary coordinate. The product $z\overline z=3^2+4^2=25$ is real, so $z^{-1}=\overline z/25=3/25-(4/25)i$.}}

\newcommand{\UnionIllustration}{%
\begin{tikzpicture}[x=1.7cm,y=1.7cm,>=Stealth,font=\small]
\draw[blue,->] (-.5,0)--(2.35,0) node[right] {$U$};
\draw[red,->] (0,-.4)--(0,1.85) node[above] {$W$};
\draw[blue,thick,->] (0,0)--(1,0) node[midway,below] {$(1,0)$};
\draw[red,thick,->] (1,0)--(1,1) node[midway,right] {$(0,1)$};
\draw[->] (0,0)--(1,1);
\draw[dashed,gray] (0,1)--(1,1);
\fill (1,1) circle (1.6pt) node[above right] {$(1,1)\notin U\cup W$};
\node[below left] at (0,0) {$0$};
\end{tikzpicture}
\figcap{The union contains only the two axes, so it omits the sum $(1,1)$. The subspace sum contains the whole plane: $(x,y)=(x,0)+(0,y)$.}}

\newcommand{\DirectSumIllustration}{%
\begin{tikzpicture}[x=.88cm,y=.88cm,>=Stealth,font=\small]
\draw[->,gray] (-.9,0)--(5.7,0) node[right,text=black] {$x$};
\draw[->,gray] (0,-1.8)--(0,4.45) node[above,text=black] {$y$};
\draw[blue] (-.8,-.8)--(4.1,4.1) node[above right] {$U:y=x$};
\draw[red] (-.8,.8)--(1.65,-1.65) node[below right] {$W:y=-x$};
\draw[blue,thick,->] (0,0)--(3,3) node[midway,above left] {$u$};
\draw[red,thick,->] (0,0)--(1,-1) node[midway,below left] {$w$};
\draw[red,thick,->] (3,3)--(4,2) node[midway,above right] {$w$};
\draw[->] (0,0)--(4,2) node[midway,below right] {$v$};
\draw[blue!50,dashed] (1,-1)--(4,2);
\fill[blue] (3,3) circle (1.6pt) node[above left] {$(3,3)$};
\fill (4,2) circle (1.6pt) node[right] {$(4,2)$};
\fill[red] (1,-1) circle (1.6pt) node[right] {$(1,-1)$};
\node[below left] at (0,0) {$0$};
\end{tikzpicture}
\figcap{$(4,2)=(3,3)+(1,-1)$, with one summand on each line. The two red arrows are the same vector. The summands are unique because $U\cap W=\{0\}$ (Proposition 7).}}

\newcommand{\ThreeLineIllustration}{%
\begin{tikzpicture}[x=1.75cm,y=1.75cm,>=Stealth,font=\small]
\draw[blue,->] (-.35,0)--(1.65,0) node[right] {$U_1$};
\draw[red,->] (0,-.3)--(0,1.55) node[above] {$U_2$};
\draw[green!50!black] (-.3,-.3)--(1.35,1.35) node[above right] {$U_3$};
\draw[blue,thick,->] (0,0)--(1,0) node[midway,below] {$(1,0)$};
\draw[red,thick,->] (1,0)--(1,1) node[midway,right] {$(0,1)$};
\draw[green!50!black,thick,->] (0,0)--(1,1);
\fill (1,1) circle (1.6pt) node[above left] {$(1,1)$};
\node[below left] at (0,0) {$0$};
\end{tikzpicture}
\figcap{The same vector is a sum of the two axis vectors, or a single vector from $U_3$:
$(1,1)=(1,0)+(0,1)+(0,0)=(0,0)+(0,0)+(1,1)$.
Pairwise zero intersections therefore do not guarantee a direct sum.}}
